\documentclass[11pt]{article}

\usepackage[letterpaper, margin=1in]{geometry}
\usepackage{graphicx}
\usepackage{enumitem}
\usepackage{hyperref}
\usepackage{xcolor}
\usepackage{titlesec}
\usepackage{booktabs}
\usepackage{parskip}
\usepackage{tabularx}

\hypersetup{
    colorlinks=true,
    linkcolor=blue,
    urlcolor=blue
}

% Helper for rough-draft notes and questions
\newcommand{\todo}[1]{\textcolor{red}{\textbf{[TODO: #1]}}}
\newcommand{\note}[1]{\textcolor{blue}{\textit{[NOTE: #1]}}}
\newcommand{\question}[1]{\textcolor{magenta}{\textbf{[QUESTION: #1]}}}

\titleformat{\section}{\large\bfseries}{\thesection}{1em}{}
\titleformat{\subsection}{\normalsize\bfseries}{\thesubsection}{1em}{}

\title{
    \vspace{-1in}
    \textbf{Universal Plastic Shredder V2.0}\\
    \large Electrical Control and Power System\\
    \normalsize \textbf{Final Report - Rough Draft}
}
\author{
    Team 07: I'm So Shredded INC\\
    Bao Nguyen, Fearghus Tyler, Yaqoub Rabiah, Fox Kang\\
    \texttt{baon@pdx.edu, fearghus@pdx.edu, yrabiah@pdx.edu, foxkang@pdx.edu}\\[4pt]
    Industry Sponsor: MME / Ichor Systems (Bridget Lannigan)\\
    Faculty Advisors: Prof. Andrew Greenberg, Dr. Jonathan Bird (EE), Robert Paxton (ME)\\[4pt]
    ECE 412 Senior Project Development I/II\\
    Portland State University, Winter--Spring 2026
}
\date{Rough Draft Submission: May 31, 2026}

\begin{document}
\maketitle
\thispagestyle{empty}

\noindent\rule{\textwidth}{0.4pt}
\note{This is a rough draft. Sections marked TODO, NOTE, or QUESTION indicate areas where we need to expand content, gather data, or get advisor input before the final submission. Bullet points are placeholders for fuller prose in the final version.}
\noindent\rule{\textwidth}{0.4pt}

\tableofcontents
\newpage

\section{Executive Summary}

\begin{itemize}
    \item Ichor Systems generates significant plastic waste (PVC, polypropylene, LDPE) during semiconductor equipment manufacturing, currently sent to landfill.
    \item Commercial industrial shredders cost \$10,000+ and consumer units cannot handle the thick rigid sheets (0.25--0.5 in).
    \item Our project designs the electrical control and power system for a custom industrial-grade plastic shredder, capable of shredding into 1-inch pieces.
    \item Key elements: VFD-controlled 3 HP motor at 70--90 RPM output, two-hand safety control, lid interlock, E-stop, jam-detection auto-recovery with 3-strike lockout, NEMA-enclosed PLC and HMI.
    \item Total budget: \$3,000 (EE portion approximately \$1,000).
\end{itemize}

\question{Should the executive summary include high-level results / status (e.g., ``shredder validated to X HP load with Y throughput rate'') once we complete testing? Currently we have signal-chain validation but not full mechanical load testing.}

\section{Introduction}

\subsection{Background and Motivation}
\begin{itemize}
    \item Ichor Systems is a Portland-area semiconductor equipment manufacturer.
    \item Their manufacturing processes produce substantial volumes of rigid plastic sheet waste.
    \item Existing disposal route: landfill -- environmentally and economically suboptimal.
    \item A previous 2023--2024 PSU capstone team built a 2 HP household shredder for the EPL. This project scales to industrial-grade.
\end{itemize}

\subsection{Project Scope}
\begin{itemize}
    \item EE team responsibility: motor control, safety circuits, PLC logic, HMI, power distribution.
    \item ME team (separate) handles mechanical frame, blade assembly, gearbox.
    \item Integration plan: mechanical structure delivered with motor mount; EE team commissions controls.
\end{itemize}

\todo{Add a one-paragraph background on dynamic recycling / circular economy framing for the introduction.}

\section{Requirements}

\subsection{Must-Have}
\begin{itemize}
    \item Shred 0.25--0.5 in plastic sheet into approximately 1-inch pieces.
    \item Controllable VFD supporting 3--5 HP motor at 70--90 RPM output shaft speed.
    \item Follow NFPA 70 (NEC) electrical design guidance.
    \item Include emergency stop, overcurrent/overload protection, at least one safety interlock.
    \item Include an HMI for operator control and status display.
    \item Stay within \$3,000 total project budget.
\end{itemize}

\subsection{Should-Have}
\begin{itemize}
    \item Clear status and fault indicator lights.
    \item Follow NEMA enclosure guidelines.
    \item Clearly labeled controls.
\end{itemize}

\subsection{May-Have (if time/budget allow)}
\begin{itemize}
    \item Wireless monitoring (read-only).
    \item Data logging of operating hours and faults.
    \item Additional safety sensing (capacitive touch, thermal).
\end{itemize}

\section{System Architecture}

\subsection{Hardware Subsystems}
\begin{itemize}
    \item \textbf{AC power input:} 120 VAC mains $\rightarrow$ main disconnect $\rightarrow$ branch breakers $\rightarrow$ VFD and 24V PSU.
    \item \textbf{VFD:} Huanyang HY02D211B-T (2.2 kW, 110V single-phase input, 220V 3-phase output boost-type). Controlled via FOR/REV/RST digital inputs.
    \item \textbf{Motor:} GE 5KE182BC205B, 3 HP, 230V (low-voltage connection), 1750 RPM, FLA 7.6 A.
    \item \textbf{Brake resistor:} 75 $\Omega$, 500 W, connected to VFD's P/Pr terminals to handle regenerative energy during deceleration and reversal.
    \item \textbf{PLC:} AutomationDirect H2-DM1E (Do-more CPU) on DL205 base, D2-08ND3 discrete input module.
    \item \textbf{HMI:} C-more Micro-Graphic EA1-T4CL 4-inch color touch panel.
    \item \textbf{Safety circuit:} Hardwired E-stop loop, lid interlock, two-hand deadman start. Motor contactor energized through safety chain.
    \item \textbf{24V DC supply:} Mean Well NDR-480-24 (480 W) for control voltage.
\end{itemize}

\subsection{Signal Flow}
\begin{itemize}
    \item Operator inputs (deadman buttons, E-stop, reset, lid interlock) $\rightarrow$ PLC discrete inputs (D2-08ND3, sourcing mode).
    \item PLC outputs (D2-08TR relay module, common to PSU --V) $\rightarrow$ VFD FOR (Y7), REV (Y6), RST (Y5) inputs in sinking (NPN) mode.
    \item VFD fault relay (FC + FA, NO active-high) $\rightarrow$ PLC X5 input. X5 = HIGH when VFD trips on fault, LOW when healthy.
    \item HMI $\leftrightarrow$ PLC via Ethernet (H2-DM1E built-in Ethernet port).
\end{itemize}

\todo{Insert system block diagram here. Convert one of the AutoCAD drawings or draw fresh in draw.io.}

\section{Hardware Design}

\subsection{Power Distribution}
\begin{itemize}
    \item Main disconnect: \todo{specify rating}.
    \item VFD branch breaker: 30 A C-curve DIN-rail (handles inrush, sized per NEC 430.52).
    \item 24V PSU branch breaker: 10 A C-curve DIN-rail.
    \item All branches fused individually at terminal blocks for downstream loads.
\end{itemize}

\subsection{Motor and VFD}
\begin{itemize}
    \item Selected Huanyang HY02D211B-T over alternatives via decision matrix (price, documentation, reliability).
    \item Programmed motor parameters: \texttt{PD141 = 230 V}, \texttt{PD142 = 7.6 A}, \texttt{PD143 = 4 poles}, \texttt{PD144 = 1750 RPM}.
    \item Over-torque detection programmed: \texttt{PD052 = 02} (fault indication on relay), \texttt{PD123 = 3}, \texttt{PD124 = 150} (\%), \texttt{PD125 = 3.0 s}.
    \item Brake resistor connected to P/Pr terminals to handle regenerative energy from rapid reversal during jam-clear cycles.
\end{itemize}

\note{The over-torque relay function (\texttt{PD052 = 12}) did not fire the relay in our specific Huanyang revision when \texttt{PD123 = 2} (continue running). Documented as a firmware-specific limitation in \texttt{vfd/Overcurrent\_Signal\_Wiring.md}. Working configuration uses \texttt{PD052 = 02} (fault indication) and \texttt{PD123 = 3} (over-torque triggers self-trip, which fires the fault relay). The PLC then runs the unjam routine including a RST pulse to clear the VFD's latched fault state between retry attempts.}

\subsection{Safety Circuit}
\begin{itemize}
    \item E-stop: NC contact pair in hardwired loop. Drops motor contactor on activation.
    \item Lid interlock: NC switch in same hardwired safety loop.
    \item Two-hand deadman: both buttons must be held continuously to enable motor.
    \item Motor contactor: K1 main contactor energized through safety chain. PLC outputs RUN command \emph{through} the contactor, so safety chain has hardware precedence over PLC.
\end{itemize}

\todo{Add safety circuit schematic. Reference NFPA 79 section 9 for safety-related stop functions.}

\subsection{Jam Detection Signal Chain}
\begin{itemize}
    \item VFD's FA/FB/FC relay programmed for Fault Indication (\texttt{PD052 = 02}).
    \item Wired as NO active-high pair: FC (relay common) connected to +24V through a 500 mA fast-blow fuse; FA (NO contact) connected to PLC X5 input on D2-08ND3 (Slot 2).
    \item PLC sees X5 = LOW during normal operation, X5 = HIGH when VFD trips on any fault (over-torque, overvoltage, overheat, etc.).
    \item PLC ladder reads \texttt{[ X5 ]} on rising edge to detect new fault events.
    \item Wire-break failure mode is handled at the safety-circuit level by the redundant hardwired E-stop loop and motor contactor, rather than at the jam-signal level.
    \item Detailed wiring in repository: \texttt{vfd/Overcurrent\_Signal\_Wiring.md}.
\end{itemize}

\section{Software Design}

\subsection{PLC Ladder Logic Architecture}
\begin{itemize}
    \item Do-more Designer stage program (similar to grafcet / SFC).
    \item States: IDLE, RUNNING, JAM\_DETECTED, RST\_PULSE, REVERSE, RESUME, LOCKOUT.
    \item Transitions driven by deadman input, X5 jam signal, jam counter, and reset button.
\end{itemize}

\subsection{State Machine}
\begin{itemize}
    \item \textbf{IDLE} $\rightarrow$ RUNNING when both deadman buttons held AND not locked out AND startup mask done.
    \item \textbf{RUNNING} $\rightarrow$ JAM\_DETECTED on rising edge of X5 (VFD trip → fault relay fires → X5 goes HIGH).
    \item \textbf{JAM\_DETECTED} $\rightarrow$ RST\_PULSE if jam count $< 3$, $\rightarrow$ LOCKOUT if jam count $\geq 3$.
    \item \textbf{RST\_PULSE} (500 ms) $\rightarrow$ REVERSE (2 s).
    \item \textbf{REVERSE} $\rightarrow$ RESUME (brief delay) $\rightarrow$ back to RUNNING.
    \item \textbf{LOCKOUT} $\rightarrow$ IDLE only via dedicated reset button.
\end{itemize}

\subsection{Jam Counter Reset Logic}
\begin{itemize}
    \item Counter increments on each X5 rising edge (each new VFD fault triggers a jam event).
    \item Resets to 0 if motor runs cleanly for 30 continuous seconds in the RUNNING state.
    \item Also resets if deadman is released for 30 continuous seconds (session reset).
    \item Whichever condition fires first wins.
\end{itemize}

\subsection{Startup Mask}
\begin{itemize}
    \item For first 5 seconds after PLC power-up, X5 is ignored.
    \item In NO active-high wiring, this guards against transient X5 pulses during VFD initialization. Less critical than under NC wiring but kept as a defensive measure.
\end{itemize}

\subsection{HMI Programming}
\begin{itemize}
    \item Status displays: Ready, Running, Jam in progress, Lockout, Fault.
    \item Operator inputs: Speed setpoint (currently fixed in PLC), Reset button.
    \item Data displayed: Motor current (read from VFD over Modbus -- \todo{verify Modbus setup}), Jam count this session, Total runtime.
\end{itemize}

\question{We have not finalized the HMI screen layout. Should we include mock-ups in the rough draft, or wait for the final?}

\section{Safety Design}

\subsection{NFPA 79 Compliance}
\begin{itemize}
    \item Wire color coding follows NFPA 79 §13.2: Black for AC line/load, Red for AC control, Blue for DC control, Blue/W stripe for DC return, Green/Yellow for ground.
    \item Documented in \texttt{electrical/Wire\_Color\_Standard.md}.
    \item All control wires labeled at both ends.
    \item All control wires terminated with crimped ferrules at terminal blocks.
\end{itemize}

\subsection{Hardwired Safety Functions}
\begin{itemize}
    \item E-stop: Category 0 stop (immediate removal of power) per ISO 13849. NC fail-safe wiring.
    \item Lid interlock: NC fail-safe; opens motor contactor when lid raised.
    \item Two-hand control: per ISO 13851 (Type IIIA at minimum) -- buttons must be pressed within 0.5 s of each other and held continuously.
    \item Overload: motor thermal protection in VFD plus separate motor contactor with overload relay.
\end{itemize}

\todo{Confirm we are meeting Category 1 or higher per ISO 13849-1. Need to walk through the SIL/PL calculation or get a hand-wave from the advisor that capstone scope doesn't require formal SIL.}

\subsection{Failure Mode Analysis}
\begin{itemize}
    \item Wire break in jam signal: PLC sees X5 stuck LOW. Jam signal cannot fire, but the hardwired E-stop loop and motor contactor provide independent safety. Operator can also see motor stop and trigger investigation via fault-light absence.
    \item VFD loses power: same as above -- relay returns to idle state which on FB-FC means contact opens, X5 LOW, locked out.
    \item PLC fails: motor contactor drops out (safety chain bypasses PLC). Motor coasts to stop.
    \item Brake resistor disconnected: VFD will trip on overvoltage during deceleration. Motor coasts to stop. Diagnosable from VFD display.
\end{itemize}

\section{Testing and Validation}

\subsection{Signal Chain Tests}
\begin{itemize}
    \item Manual jumper test (FA-FC short): confirmed PLC X5 reads LOW correctly -- wiring intact.
    \item Guaranteed-trip test (\texttt{PD124 = 1}, \texttt{PD125 = 5.0}): VFD trips reliably on motor run, fault relay fires, PLC X5 drops.
    \item Threshold tuning: 150\% (\texttt{PD124 = 150}) chosen for production based on no-load current (0.5 A) versus expected jam current (above 11 A actual).
\end{itemize}

\subsection{Load Testing}
\todo{Schedule load testing with actual plastic samples once mechanical integration complete. Need ME team coordination.}

\subsection{Safety Validation}
\begin{itemize}
    \item E-stop response time: \todo{measure with scope}.
    \item Lid interlock response: \todo{measure}.
    \item Two-hand control timing: \todo{verify both buttons must be held continuously, simultaneous release stops within X ms}.
\end{itemize}

\section{Results}

\subsection{What Works}
\begin{itemize}
    \item All motor parameters programmed and verified.
    \item Over-torque detection chain validated (VFD detects, relay fires, PLC reads, ladder reacts).
    \item Fail-safe NC wiring confirmed -- system correctly enters fault state on signal chain failure.
    \item HMI communicates with PLC over Ethernet.
\end{itemize}

\subsection{Open Issues}
\begin{itemize}
    \item Vendor-specific Huanyang firmware quirk: \texttt{PD052 = 12} does not fire relay in continue-running mode. Working around with \texttt{PD052 = 02} + \texttt{PD123 = 3}.
    \item PLC ladder logic not fully complete -- state machine implemented, but unjam routine timing needs tuning with real loaded motor.
    \item Mechanical integration pending -- have not yet tested under real shredding load.
\end{itemize}

\question{How much detail should we include about the Huanyang firmware quirk? It is an interesting design lesson (vendor lock-in, never trust the datasheet completely) but distracts from the high-level results.}

\section{Discussion}

\subsection{Design Tradeoffs}
\begin{itemize}
    \item Cost versus reliability: chose Huanyang VFD (\$150) over ATO (\$450). Documented quirks in repository for future maintenance.
    \item 110V single-phase versus 220V split-phase mains: 110V chosen for receptacle availability in PSU lab; trade-off is double the input current and need for a dedicated 30 A circuit.
    \item NC fail-safe versus NO active-high jam signal: chose NC for safety; trade-off is potential nuisance lockouts from broken wires (acceptable for shredder).
\end{itemize}

\subsection{Lessons Learned}
\begin{itemize}
    \item Always verify VFD relay function with multimeter before assuming firmware works as documented.
    \item Document deviations from datasheet behavior in version control so future teams can find them.
    \item Power cycle after parameter changes -- some VFDs require it for save to fully take effect.
    \item Use ferrules on all control wires -- loose strand connections caused multiple debugging hours during commissioning.
\end{itemize}

\subsection{Project Management Notes}
\todo{Add a paragraph on how the team coordinated, what tools we used (GitHub, weekly meetings, etc.), and how schedule slipped or held.}

\section{Future Work}

\begin{itemize}
    \item Wireless monitoring (read-only) -- could expose PLC data via the H2-DM1E's built-in web server or an ESP32 bridge.
    \item Data logging: total runtime, jam events per shift, faults. Useful for predictive maintenance.
    \item Predictive jam detection: train a small model on current waveform shape just before jams to detect them earlier.
    \item Mechanical integration with feed hopper that auto-meters plastic into the blade.
    \item Higher-capacity model: scale to 5--10 HP for higher throughput.
\end{itemize}

\section{Ethics and Design Considerations}

See \texttt{docs/Ethics\_Design\_Evaluation.docx} for the full filled-out form. Summary:

\begin{itemize}
    \item \textbf{Public health and safety:} Yes -- multiple redundant safety interlocks, NFPA 70 / NEMA compliance, hardwired safety chain.
    \item \textbf{Environmental:} Yes -- core project purpose is diverting plastic from landfill.
    \item \textbf{Economic:} Yes -- significant cost savings versus commercial alternatives, careful component selection within budget.
    \item \textbf{Global:} Possibly -- approach is replicable globally, components internationally sourced.
    \item \textbf{Cultural:} No -- industrial machine for B2B use, no culture-specific design.
\end{itemize}

\section{References}

\begin{itemize}
    \item NFPA 70 (2023): National Electrical Code.
    \item NFPA 79 (2024): Electrical Standard for Industrial Machinery.
    \item ISO 13849-1:2023: Safety-related parts of control systems.
    \item ISO 13851:2019: Two-hand control devices.
    \item Huanyang HY-series VFD User Manual.
    \item AutomationDirect DL205 / Do-more User Manual.
    \item GE 5KE Frame 182T Motor Datasheet.
\end{itemize}

\todo{Add full reference list with URLs / publication info before final.}

\section{Appendices}

\subsection{Appendix A: Bill of Materials}
\todo{Insert latest BOM (\texttt{BOM/BOM\_02\_20\_26.xlsx}) as a formatted table or summary.}

\subsection{Appendix B: PLC I/O Point List}
\todo{Insert from \texttt{PLC/Point List.pdf}.}

\subsection{Appendix C: Wiring Diagrams}
\todo{Insert exported PDFs from AutoCAD: power, VFD/motor, PLC/HMI, power distribution.}

\subsection{Appendix D: VFD Parameter Settings}
See repository: \texttt{vfd/Overcurrent\_Signal\_Wiring.md}, \texttt{vfd/Brake\_Resistor\_Wiring.md}, \texttt{motor/Motor\_Nameplate.md}.

\subsection{Appendix E: Test Plan}
See \texttt{docs/Test\_Plan\_v1.0.md}.

\vfill
\noindent\rule{\textwidth}{0.4pt}
\noindent\textit{End of rough draft. Comments and questions welcome -- please use Google Docs share for inline review (forthcoming).}
\noindent\rule{\textwidth}{0.4pt}

\end{document}
